\selectlanguage{english}%
\noindent \begin{center}
\section*{Abstract}
\par\end{center}

ROS provides a set of tools and libraries for developing robotic applications
that facilitate communication between different system components, as well as the integration of sensors and actuators.
There are various applications based on this technology that allow for the visualization and control of robots; however,
most are focused on desktop environments.
An example of this is the well-known \textit{RViz} tool, which allows for the visualization of sensor data and the simulation of robots,
but requires a personal computer to run,
which limits users' ability to interact with them remotely and in real time, especially in situations
where mobility and flexibility are required.
There are other examples that run on mobile platforms, such as \textit{ROS-Mobile} or \textit{Robot Steering App}, but,
again, these are menu-based solutions and do not allow for direct interaction with the physical environment, which hinders
understanding and effective control.

Augmented reality, on the other hand, allows interfaces to be built in such a way that the real world and the virtual world are integrated into one through
a device capable of capturing images. Data using this paradigm will be displayed to the user via
three-dimensional or textual objects superimposed on the real world, which allows for
the creation of simpler and more manageable representations for users.

Therefore, we propose the development of a cross-platform augmented reality mobile app for visualizing
robot data using ROS2 and Unity,
so that users can access robot information from any device, including smartphones and tablets,
thereby reducing the barrier between the physical and digital worlds and enabling more intuitive and efficient interaction.

\vspace{0.5cm}
\begin{flushleft}
\textbf{Keywords}: ROS2, Unity, smartphone, visualization, control, augmented reality
\end{flushleft}
\selectlanguage{spanish}%